% =====================================================================
%  แม่แบบวิทยานิพนธ์ภาษาไทย (XeLaTeX)
%  คอมไพล์:  latexmk   (ตั้งค่า XeLaTeX + Biber ไว้ใน latexmkrc แล้ว)
%  ผลลัพธ์:  build/main.pdf  (ไฟล์ชั่วคราวทั้งหมดอยู่ใน build/)
% =====================================================================
\documentclass{thaithesis}
% ตัวเลือกเพิ่มเติม เช่น
%   \documentclass[pagebottom]{thaithesis}          เลขหน้ากึ่งกลางด้านล่าง
%   \documentclass[chapterpagenumber]{thaithesis}   แสดงเลขหน้าในหน้าแรกของบท
%   \documentclass[systemfont]{thaithesis}          ข้ามไฟล์ใน fonts/ ใช้ฟอนต์ที่ติดตั้งในเครื่อง

% ขอบกระดาษตามข้อกำหนดของสถาบัน (ค่าเริ่มต้น: บน 1.5" ซ้าย 1.5" ขวา 1" ล่าง 1")
% \geometry{top=3.5cm, left=3.5cm, right=2.5cm, bottom=3.5cm}

% ระยะย่อหน้า (ค่าเริ่มต้น 0.5 นิ้ว)
% \setlength\thesisindent{1.5cm}\setlength\parindent{\thesisindent}

% ---------------------------------------------------------------------
%  ข้อมูลวิทยานิพนธ์  (แก้ไขเป็นข้อมูลจริง)
% ---------------------------------------------------------------------
\thesistypeth{วิทยานิพนธ์}      % หรือ ปริญญานิพนธ์ / สารนิพนธ์ / โครงงาน
\thesistypeen{Thesis}           % หรือ Dissertation / Independent Study / Project

% ใช้ \\ เพื่อขึ้นบรรทัดใหม่ให้ชื่อเรื่องเป็นรูปสามเหลี่ยมกลับหัว
\titleth{การพัฒนาระบบผู้ช่วยสอนด้วยปัญญาประดิษฐ์เชิงสร้างสรรค์\\สำหรับรายวิชาการเขียนโปรแกรมเบื้องต้น}
\titleen{Development of a Generative AI Teaching Assistant\\for an Introductory Programming Course}

\authorth{นายสมชาย ใจดี}
\authoren{Mr. Somchai Jaidee}

\degreeth{วิศวกรรมศาสตรมหาบัณฑิต}
\degreeen{Master of Engineering}
\programth{วิศวกรรมคอมพิวเตอร์}
\programen{Computer Engineering}
\facultyth{คณะวิศวกรรมศาสตร์}
\facultyen{Faculty of Engineering}
\instituteth{มหาวิทยาลัยตัวอย่าง}
\instituteen{Example University}
\academicyear{2569}{2026}       % {พ.ศ.}{ค.ศ.}

\advisorth{ผู้ช่วยศาสตราจารย์ ดร.สมหญิง ตั้งใจเรียน}
\advisoren{Asst. Prof. Dr. Somying Tangjairian}
\coadvisorth{}                  % เว้นว่างได้ถ้าไม่มี
\coadvisoren{}

\deanth{รองศาสตราจารย์ ดร.ใจดี มีเมตตา}
\committee{ประธานกรรมการ}{รองศาสตราจารย์ ดร.สมศักดิ์ รักวิชา}
\committee{กรรมการ}{ดร.วิภาวี ศรีสว่าง}
\committee{กรรมการและอาจารย์ที่ปรึกษา}{ผู้ช่วยศาสตราจารย์ ดร.สมหญิง ตั้งใจเรียน}

\keywordsth{ปัญญาประดิษฐ์เชิงสร้างสรรค์, แบบจำลองภาษาขนาดใหญ่, ผู้ช่วยสอน, การเขียนโปรแกรม}
\keywordsen{Generative AI, Large Language Models, Teaching Assistant, Programming Education}

% ตราสถาบัน: วางไฟล์ figures/logo.png (หรือ .pdf/.jpg) แล้วระบุชื่อไฟล์โดยไม่ต้องใส่นามสกุล
\logo[3.5cm]{figures/logo}

% ---------------------------------------------------------------------
%  บรรณานุกรม (Biber)
%  - แบบ IEEE (ตัวเลข [1]):  style=ieee
%  - แบบ APA (ผู้แต่ง, ปี):   style=apa
% ---------------------------------------------------------------------
\usepackage[backend=biber, style=ieee, sorting=none]{biblatex}
\addbibresource{references.bib}

% ---------------------------------------------------------------------
%  แพ็กเกจเพิ่มเติมสำหรับตัวอย่าง
% ---------------------------------------------------------------------
\usepackage{tikz}
\usetikzlibrary{positioning, arrows.meta}
\usepackage{pgfplots}
\pgfplotsset{compat=1.18}
\usepackage{listings}
\lstset{
  basicstyle=\ttfamily\small,
  columns=fullflexible, keepspaces=true,
  frame=single, framerule=0.4pt, rulecolor=\color{black!40},
  numbers=left, numberstyle=\ttfamily\scriptsize\color{black!50},
  keywordstyle=\bfseries\color{blue!60!black},
  commentstyle=\itshape\color{green!40!black},
  stringstyle=\color{red!60!black},
  breaklines=true, showstringspaces=false,
  xleftmargin=2em, framexleftmargin=1.5em}

% =====================================================================
\begin{document}

% ---------- ส่วนปก (ไม่แสดงเลขหน้า) ----------
\makecover        % ปกนอก
\maketitleth      % ปกในภาษาไทย
\maketitleen      % ปกในภาษาอังกฤษ
\makeapproval     % หน้าอนุมัติ

% ---------- ส่วนนำ (เลขหน้า ก ข ค ...) ----------
\frontmatter
% =====================================================================
%  บทคัดย่อภาษาไทย
%  ข้อมูลหัวข้อ ชื่อผู้เขียน อาจารย์ที่ปรึกษา ฯลฯ ดึงมาจาก main.tex อัตโนมัติ
%  คำสำคัญ กำหนดด้วย \keywordsth{...} ใน main.tex
% =====================================================================
\begin{abstractth}
งานวิจัยนี้มีวัตถุประสงค์เพื่อพัฒนาระบบผู้ช่วยสอนที่ใช้ปัญญาประดิษฐ์เชิงสร้างสรรค์ (Generative AI) สำหรับรายวิชาการเขียนโปรแกรมเบื้องต้น และเพื่อประเมินประสิทธิภาพของระบบในการช่วยให้ผู้เรียนเข้าใจแนวคิดการเขียนโปรแกรมและแก้ไขข้อผิดพลาดของโค้ดได้ด้วยตนเอง ระบบที่พัฒนาขึ้นใช้แบบจำลองภาษาขนาดใหญ่ (Large Language Model: LLM) ร่วมกับเทคนิคการค้นคืนข้อมูลเสริมการสร้าง (Retrieval-Augmented Generation: RAG) เพื่ออ้างอิงเนื้อหาจากเอกสารประกอบการสอนของรายวิชา

กลุ่มตัวอย่างคือนักศึกษาชั้นปีที่ 1 จำนวน 60 คน แบ่งเป็นกลุ่มทดลองและกลุ่มควบคุมกลุ่มละ 30 คน ผลการวิจัยพบว่า กลุ่มทดลองมีคะแนนผลสัมฤทธิ์ทางการเรียนหลังเรียนสูงกว่ากลุ่มควบคุมอย่างมีนัยสำคัญทางสถิติที่ระดับ .05 และมีความพึงพอใจต่อระบบในระดับมากที่สุด (ค่าเฉลี่ย 4.52 จากคะแนนเต็ม 5)

\textit{(ข้อความนี้เป็นตัวอย่าง บทคัดย่อควรมีความยาวไม่เกิน 1 หน้า ประกอบด้วยวัตถุประสงค์ วิธีการ ผลการวิจัย และข้อสรุปโดยย่อ)}
\end{abstractth}

% =====================================================================
%  บทคัดย่อภาษาอังกฤษ (ABSTRACT)
%  คำสำคัญ กำหนดด้วย \keywordsen{...} ใน main.tex
% =====================================================================
\begin{abstracten}
This research aims to develop a teaching assistant system powered by Generative AI for an introductory programming course and to evaluate its effectiveness in helping students understand programming concepts and debug their own code. The system combines a Large Language Model (LLM) with Retrieval-Augmented Generation (RAG) so that its answers are grounded in the course materials.

The sample consisted of 60 first-year students, divided into an experimental group and a control group of 30 students each. The results showed that the post-test achievement scores of the experimental group were significantly higher than those of the control group at the .05 level, and that students were highly satisfied with the system (mean 4.52 out of 5).

\textit{(This is sample text. The abstract should not exceed one page.)}
\end{abstracten}

% =====================================================================
%  กิตติกรรมประกาศ
%  ชื่อผู้เขียนจะแสดงชิดขวาท้ายหน้าโดยอัตโนมัติ
% =====================================================================
\begin{acknowledgements}
วิทยานิพนธ์ฉบับนี้สำเร็จลุล่วงได้ด้วยความกรุณาอย่างยิ่งจากผู้ช่วยศาสตราจารย์ ดร.สมหญิง ตั้งใจเรียน อาจารย์ที่ปรึกษา ที่ได้ให้คำแนะนำ ข้อคิดเห็น และตรวจแก้ไขข้อบกพร่องต่าง ๆ ด้วยความเอาใจใส่ตลอดระยะเวลาของการทำวิจัย ผู้เขียนขอกราบขอบพระคุณเป็นอย่างสูงไว้ ณ โอกาสนี้

ขอขอบพระคุณคณะกรรมการสอบวิทยานิพนธ์ทุกท่านที่ได้กรุณาให้ข้อเสนอแนะอันเป็นประโยชน์ ขอบคุณนักศึกษาทุกคนที่ให้ความร่วมมือในการเก็บรวบรวมข้อมูล และขอบคุณมหาวิทยาลัยที่สนับสนุนทุนการศึกษาและสิ่งอำนวยความสะดวกในการวิจัย

ท้ายที่สุดนี้ ผู้เขียนขอกราบขอบพระคุณบิดา มารดา และครอบครัว ที่ให้กำลังใจและสนับสนุนผู้เขียนเสมอมา

% เครดิตแม่แบบ -- ขอความกรุณาคงประโยคนี้ไว้เพื่อให้เครดิตผู้พัฒนาแม่แบบ ขอบคุณครับ
% Template credit -- please keep this sentence to acknowledge the template author. Thank you!
รูปแบบของวิทยานิพนธ์ฉบับนี้จัดทำด้วยแม่แบบ LaTeX ของ ดร.ชยันต์ คงทองวัฒนา~\cite{kongtongvattana2026thaithesis}
\end{acknowledgements}

\tableofcontents
\listoftables
\listoffigures
% =====================================================================
%  คำอธิบายสัญลักษณ์และคำย่อ
%  รูปแบบ: \abbr{สัญลักษณ์หรือคำย่อ}{ความหมาย}
% =====================================================================
\begin{abbreviations}
\abbr{AI}{ปัญญาประดิษฐ์ (Artificial Intelligence)}
\abbr{GAN}{เครือข่ายปฏิปักษ์เชิงสร้าง (Generative Adversarial Network)}
\abbr{LLM}{แบบจำลองภาษาขนาดใหญ่ (Large Language Model)}
\abbr{RAG}{การค้นคืนข้อมูลเสริมการสร้าง (Retrieval-Augmented Generation)}
\abbr{VAE}{ตัวเข้ารหัสอัตโนมัติแบบแปรผัน (Variational Autoencoder)}
\abbr{$\mathbf{x}$}{เวกเตอร์ข้อมูลนำเข้า}
\abbr{$\theta$}{พารามิเตอร์ของแบบจำลอง}
\abbr{$\mathcal{L}$}{ฟังก์ชันค่าสูญเสีย (loss function)}
\end{abbreviations}


% ---------- เนื้อหา (เลขหน้า 1 2 3 ...) ----------
\mainmatter
% =====================================================================
%  บทที่ 1  บทนำ
% =====================================================================
\chapter{บทนำ}
\label{ch:intro}

\section{ความเป็นมาและความสำคัญของปัญหา}

ปัญญาประดิษฐ์เชิงสร้างสรรค์ (Generative Artificial Intelligence) เป็นเทคโนโลยีที่สามารถสร้างเนื้อหาใหม่ เช่น ข้อความ ภาพ เสียง และโค้ดโปรแกรม ได้จากข้อมูลที่ใช้ฝึกสอน ความก้าวหน้าของสถาปัตยกรรมทรานส์ฟอร์เมอร์ (Transformer)~\cite{vaswani2017attention} ทำให้เกิดแบบจำลองภาษาขนาดใหญ่ที่สามารถเรียนรู้จากตัวอย่างเพียงไม่กี่ตัวอย่างได้~\cite{brown2020language} ส่งผลให้มีการนำเทคโนโลยีนี้มาประยุกต์ใช้ในด้านการศึกษาอย่างแพร่หลาย

ในรายวิชาการเขียนโปรแกรมเบื้องต้น ผู้เรียนจำนวนมากประสบปัญหาในการทำความเข้าใจข้อความแจ้งข้อผิดพลาด (error message) และไม่สามารถแก้ไขโค้ดได้ด้วยตนเอง ขณะที่ผู้สอนมีเวลาจำกัดในการให้คำแนะนำแก่ผู้เรียนเป็นรายบุคคล\footnote{ตัวอย่างเชิงอรรถ: ข้อมูลจากการสำรวจเบื้องต้นของผู้วิจัยในภาคการศึกษาที่ 1/2568} ด้วยเหตุนี้ ผู้วิจัยจึงสนใจพัฒนาระบบผู้ช่วยสอนที่สามารถให้คำแนะนำแก่ผู้เรียนได้ตลอดเวลา

\section{วัตถุประสงค์ของการวิจัย}

\begin{enumerate}
  \item เพื่อพัฒนาระบบผู้ช่วยสอนด้วยปัญญาประดิษฐ์เชิงสร้างสรรค์สำหรับรายวิชาการเขียนโปรแกรมเบื้องต้น
  \item เพื่อเปรียบเทียบผลสัมฤทธิ์ทางการเรียนของผู้เรียนที่ใช้และไม่ใช้ระบบผู้ช่วยสอน
  \item เพื่อศึกษาความพึงพอใจของผู้เรียนที่มีต่อระบบผู้ช่วยสอน
\end{enumerate}

\section{ขอบเขตของการวิจัย}

\subsection{ขอบเขตด้านเนื้อหา}
ระบบครอบคลุมเนื้อหาการเขียนโปรแกรมภาษา Python ได้แก่ ตัวแปร เงื่อนไข การวนซ้ำ ฟังก์ชัน และรายการ (list)

\subsection{ขอบเขตด้านประชากรและกลุ่มตัวอย่าง}
ประชากรคือนักศึกษาชั้นปีที่ 1 ที่ลงทะเบียนเรียนรายวิชาการเขียนโปรแกรมเบื้องต้น ภาคการศึกษาที่ 1 ปีการศึกษา 2569 กลุ่มตัวอย่างจำนวน 60 คน ได้มาโดยการสุ่มแบบกลุ่ม (cluster random sampling)

\section{ประโยชน์ที่คาดว่าจะได้รับ}

\begin{itemize}
  \item ได้ระบบผู้ช่วยสอนที่ผู้เรียนสามารถใช้งานได้ตลอด 24 ชั่วโมง
  \item ผู้สอนมีเวลามากขึ้นในการดูแลผู้เรียนที่ต้องการความช่วยเหลือเป็นพิเศษ
  \item เป็นแนวทางในการประยุกต์ใช้ปัญญาประดิษฐ์เชิงสร้างสรรค์ในรายวิชาอื่น
\end{itemize}

\section{นิยามศัพท์เฉพาะ}

\textbf{ปัญญาประดิษฐ์เชิงสร้างสรรค์} หมายถึง ระบบปัญญาประดิษฐ์ที่สามารถสร้างข้อมูลใหม่ที่มีลักษณะคล้ายกับข้อมูลที่ใช้ฝึกสอน

\textbf{ผู้ช่วยสอน} หมายถึง ระบบซอฟต์แวร์ที่ตอบคำถามและให้คำแนะนำแก่ผู้เรียนในรายวิชา โดยอ้างอิงจากเอกสารประกอบการสอน

% =====================================================================
%  บทที่ 2  ทฤษฎีและงานวิจัยที่เกี่ยวข้อง
% =====================================================================
\chapter{ทฤษฎีและงานวิจัยที่เกี่ยวข้อง}
\label{ch:literature}

ในบทนี้กล่าวถึงทฤษฎีพื้นฐานของปัญญาประดิษฐ์เชิงสร้างสรรค์และงานวิจัยที่เกี่ยวข้อง ซึ่งเป็นพื้นฐานในการออกแบบระบบในบทที่~\ref{ch:method}

\section{แบบจำลองเชิงสร้าง}

แบบจำลองเชิงสร้าง (generative model) มีเป้าหมายเพื่อเรียนรู้การแจกแจงความน่าจะเป็นของข้อมูล $p(\mathbf{x})$ แล้วสุ่มตัวอย่างใหม่จากการแจกแจงนั้น แบบจำลองที่สำคัญมีดังนี้

\subsection{เครือข่ายปฏิปักษ์เชิงสร้าง}

เครือข่ายปฏิปักษ์เชิงสร้าง (GAN)~\cite{goodfellow2014generative} ประกอบด้วยเครือข่ายสองส่วน คือ ตัวสร้าง (generator) $G$ และตัวจำแนก (discriminator) $D$ ซึ่งแข่งขันกันตามฟังก์ชันเป้าหมายในสมการที่~\eqref{eq:gan}
\begin{equation}
  \min_{G}\max_{D}\;
  \mathbb{E}_{\mathbf{x}\sim p_{\text{data}}}\!\left[\log D(\mathbf{x})\right]
  + \mathbb{E}_{\mathbf{z}\sim p_{\mathbf{z}}}\!\left[\log\bigl(1-D(G(\mathbf{z}))\bigr)\right]
  \label{eq:gan}
\end{equation}

\subsection{ตัวเข้ารหัสอัตโนมัติแบบแปรผัน}

ตัวเข้ารหัสอัตโนมัติแบบแปรผัน (VAE)~\cite{kingma2014auto} เรียนรู้ตัวแปรแฝง $\mathbf{z}$ โดยหาค่าสูงสุดของขอบเขตล่างของหลักฐาน (evidence lower bound) ดังสมการที่~\eqref{eq:elbo}
\begin{equation}
  \mathcal{L}(\theta,\phi;\mathbf{x}) =
  \mathbb{E}_{q_\phi(\mathbf{z}\mid\mathbf{x})}\left[\log p_\theta(\mathbf{x}\mid\mathbf{z})\right]
  - D_{\mathrm{KL}}\!\left(q_\phi(\mathbf{z}\mid\mathbf{x})\,\|\,p(\mathbf{z})\right)
  \label{eq:elbo}
\end{equation}

\subsection{แบบจำลองการแพร่}

แบบจำลองการแพร่ (diffusion model)~\cite{ho2020denoising} ค่อย ๆ เติมสัญญาณรบกวนลงในข้อมูลจนกลายเป็นสัญญาณรบกวนแบบเกาส์เซียน แล้วฝึกเครือข่ายให้ย้อนกระบวนการดังกล่าว โดยใช้ฟังก์ชันค่าสูญเสียอย่างง่ายดังสมการที่~\eqref{eq:ddpm}
\begin{equation}
  \mathcal{L}_{\text{simple}} =
  \mathbb{E}_{t,\mathbf{x}_0,\boldsymbol{\epsilon}}
  \left[\bigl\lVert \boldsymbol{\epsilon}
  - \boldsymbol{\epsilon}_\theta\!\left(\sqrt{\bar{\alpha}_t}\,\mathbf{x}_0
  + \sqrt{1-\bar{\alpha}_t}\,\boldsymbol{\epsilon},\, t\right)\bigr\rVert^2\right]
  \label{eq:ddpm}
\end{equation}

\section{สถาปัตยกรรมทรานส์ฟอร์เมอร์}

หัวใจสำคัญของทรานส์ฟอร์เมอร์~\cite{vaswani2017attention} คือกลไกความสนใจ (attention) ซึ่งคำนวณได้ดังสมการที่~\eqref{eq:attention} เมื่อ $Q$, $K$ และ $V$ คือเมทริกซ์คิวรี คีย์ และค่า ตามลำดับ และ $d_k$ คือมิติของคีย์
\begin{equation}
  \operatorname{Attention}(Q,K,V) = \operatorname{softmax}\!\left(\frac{QK^{\top}}{\sqrt{d_k}}\right)V
  \label{eq:attention}
\end{equation}

\section{งานวิจัยที่เกี่ยวข้อง}

ตารางที่~\ref{tab:related} สรุปแบบจำลองเชิงสร้างที่สำคัญซึ่งเกี่ยวข้องกับงานวิจัยนี้

\begin{table}[htbp]
  \caption{สรุปแบบจำลองเชิงสร้างที่สำคัญ}
  \label{tab:related}
  \centering
  \begin{tabularx}{\linewidth}{@{}l c X@{}}
    \toprule
    \textbf{แบบจำลอง} & \textbf{ปี} & \textbf{จุดเด่น}\\
    \midrule
    GAN~\cite{goodfellow2014generative} & 2014 & สร้างภาพได้คมชัด แต่การฝึกสอนไม่เสถียร\\
    VAE~\cite{kingma2014auto} & 2014 & มีพื้นที่แฝงที่ต่อเนื่อง ตีความได้ง่าย\\
    Transformer~\cite{vaswani2017attention} & 2017 & ประมวลผลลำดับแบบขนานด้วยกลไกความสนใจ\\
    DDPM~\cite{ho2020denoising} & 2020 & คุณภาพภาพสูงและฝึกสอนได้เสถียร\\
    GPT-3~\cite{brown2020language} & 2020 & เรียนรู้จากตัวอย่างเพียงไม่กี่ตัวอย่าง (few-shot)\\
    \bottomrule
  \end{tabularx}
\end{table}

นอกจากเอกสารภาษาอังกฤษแล้ว สามารถอ้างอิงเอกสารภาษาไทยได้ในรูปแบบเดียวกัน~\cite{thai-example}

% =====================================================================
%  บทที่ 3  วิธีดำเนินการวิจัย
% =====================================================================
\chapter{วิธีดำเนินการวิจัย}
\label{ch:method}

\section{ภาพรวมของระบบ}

ระบบผู้ช่วยสอนที่พัฒนาขึ้นประกอบด้วยสามส่วนหลัก ได้แก่ ส่วนค้นคืนเอกสาร ส่วนสร้างคำตอบ และส่วนติดต่อผู้ใช้ ดังแสดงในภาพที่~\ref{fig:overview}

\begin{figure}[htbp]
  \centering
  \begin{tikzpicture}[
      node distance=8mm and 10mm,
      box/.style={draw, rounded corners=3pt, thick, align=center,
                  minimum width=30mm, minimum height=12mm, fill=blue!6},
      db/.style={draw, thick, align=center, minimum width=30mm,
                 minimum height=12mm, fill=orange!10},
      arr/.style={-{Stealth[length=2.5mm]}, thick}]
    \node[box] (q)   {คำถามจากผู้เรียน};
    \node[box, right=of q] (ret) {ส่วนค้นคืนเอกสาร\\(Retriever)};
    \node[db,  below=of ret] (docs) {เอกสารประกอบ\\การสอน};
    \node[box, right=of ret] (llm) {แบบจำลองภาษา\\(LLM)};
    \node[box, below=of llm] (ans) {คำตอบและ\\คำแนะนำ};
    \draw[arr] (q) -- (ret);
    \draw[arr] (docs) -- (ret);
    \draw[arr] (ret) -- node[above]{บริบท} (llm);
    \draw[arr] (llm) -- (ans);
  \end{tikzpicture}
  \caption{ภาพรวมของระบบผู้ช่วยสอนแบบ RAG}
  \label{fig:overview}
\end{figure}

\section{การเตรียมข้อมูล}

เอกสารประกอบการสอนทั้งหมดถูกแบ่งเป็นส่วนย่อย (chunk) ขนาดประมาณ 500 โทเคน แล้วแปลงเป็นเวกเตอร์ฝังตัว (embedding) เพื่อจัดเก็บในฐานข้อมูลเวกเตอร์ ความคล้ายคลึงระหว่างคำถาม $\mathbf{q}$ กับเอกสาร $\mathbf{d}$ คำนวณด้วยความคล้ายคลึงโคไซน์ดังสมการที่~\eqref{eq:cosine}
\begin{equation}
  \operatorname{sim}(\mathbf{q},\mathbf{d}) =
  \frac{\mathbf{q}\cdot\mathbf{d}}{\lVert\mathbf{q}\rVert\,\lVert\mathbf{d}\rVert}
  \label{eq:cosine}
\end{equation}

\section{การพัฒนาระบบ}

ตัวอย่างโค้ดสำหรับค้นคืนเอกสารที่เกี่ยวข้องแสดงดังโปรแกรมต่อไปนี้

\begin{lstlisting}[language=Python]
import numpy as np

def retrieve(query_vec, doc_vecs, docs, k=3):
    """Return the k most similar documents (cosine similarity)."""
    q = query_vec / np.linalg.norm(query_vec)
    d = doc_vecs / np.linalg.norm(doc_vecs, axis=1, keepdims=True)
    scores = d @ q
    top = np.argsort(scores)[::-1][:k]
    return [docs[i] for i in top]
\end{lstlisting}

\section{เครื่องมือที่ใช้ในการวิจัย}

\begin{enumerate}
  \item ระบบผู้ช่วยสอนที่พัฒนาขึ้น
  \item แบบทดสอบวัดผลสัมฤทธิ์ทางการเรียน จำนวน 30 ข้อ
  \item แบบสอบถามความพึงพอใจ แบบมาตราส่วนประมาณค่า 5 ระดับ (ภาคผนวก~\ref{app:questionnaire})
\end{enumerate}

\section{การวิเคราะห์ข้อมูล}

เปรียบเทียบคะแนนหลังเรียนระหว่างกลุ่มทดลองและกลุ่มควบคุมด้วยสถิติทดสอบที (independent samples $t$-test) ดังสมการที่~\eqref{eq:ttest}
\begin{equation}
  t = \frac{\bar{x}_1 - \bar{x}_2}{\sqrt{\dfrac{s_1^2}{n_1} + \dfrac{s_2^2}{n_2}}}
  \label{eq:ttest}
\end{equation}

% =====================================================================
%  บทที่ 4  ผลการวิจัย
% =====================================================================
\chapter{ผลการวิจัย}
\label{ch:results}

\section{ผลการเปรียบเทียบผลสัมฤทธิ์ทางการเรียน}

ผลการเปรียบเทียบคะแนนผลสัมฤทธิ์ทางการเรียนหลังเรียนระหว่างกลุ่มทดลองและกลุ่มควบคุมแสดงในตารางที่~\ref{tab:ttest}

\begin{table}[htbp]
  \caption{ผลการเปรียบเทียบคะแนนหลังเรียนระหว่างกลุ่มทดลองและกลุ่มควบคุม}
  \label{tab:ttest}
  \centering
  \begin{tabular}{@{}l c c c c c@{}}
    \toprule
    \textbf{กลุ่ม} & $n$ & $\bar{x}$ & $S.D.$ & $t$ & $p$\\
    \midrule
    กลุ่มทดลอง  & 30 & 24.37 & 2.81 & \multirow{2}{*}{4.12} & \multirow{2}{*}{.000*}\\
    กลุ่มควบคุม & 30 & 21.13 & 3.24 &  & \\
    \bottomrule
    \multicolumn{6}{@{}l}{\small * มีนัยสำคัญทางสถิติที่ระดับ .05}
  \end{tabular}
\end{table}

จากตารางที่~\ref{tab:ttest} พบว่ากลุ่มทดลองมีคะแนนเฉลี่ยหลังเรียนสูงกว่ากลุ่มควบคุมอย่างมีนัยสำคัญทางสถิติที่ระดับ .05 และเมื่อพิจารณาคะแนนรายหัวข้อ พบว่าหัวข้อที่มีความแตกต่างมากที่สุดคือการวนซ้ำ ดังภาพที่~\ref{fig:scores}

\begin{figure}[htbp]
  \centering
  \begin{tikzpicture}
    \begin{axis}[
        ybar, bar width=9pt, width=0.9\linewidth, height=6.5cm,
        ymin=0, ymax=10, ylabel={คะแนนเฉลี่ย},
        symbolic x coords={ตัวแปร, เงื่อนไข, การวนซ้ำ, ฟังก์ชัน, รายการ},
        xtick=data, enlarge x limits=0.15,
        legend style={at={(0.5,-0.18)}, anchor=north, legend columns=-1,
                      draw=none, /tikz/every even column/.append style={column sep=1em}},
        tick label style={font=\small}, label style={font=\small}]
      \addplot[fill=blue!55]   coordinates {(ตัวแปร,8.9) (เงื่อนไข,8.1) (การวนซ้ำ,7.6) (ฟังก์ชัน,7.2) (รายการ,7.0)};
      \addplot[fill=orange!60] coordinates {(ตัวแปร,8.5) (เงื่อนไข,7.4) (การวนซ้ำ,5.9) (ฟังก์ชัน,6.3) (รายการ,6.1)};
      \legend{กลุ่มทดลอง, กลุ่มควบคุม}
    \end{axis}
  \end{tikzpicture}
  \caption{คะแนนเฉลี่ยรายหัวข้อของกลุ่มทดลองและกลุ่มควบคุม}
  \label{fig:scores}
\end{figure}

\section{ผลการประเมินความพึงพอใจ}

ตารางที่~\ref{tab:satisfaction} แสดงตัวอย่างตารางยาวที่ต่อข้ามหน้าได้ (\texttt{longtable}) โดยหัวตารางจะแสดงซ้ำพร้อมคำว่า ``(ต่อ)'' ในหน้าถัดไปโดยอัตโนมัติ

\begin{longtable}{@{}p{0.62\linewidth} c c l@{}}
  \caption{ความพึงพอใจของผู้เรียนที่มีต่อระบบผู้ช่วยสอน}
  \label{tab:satisfaction}\\
  \toprule
  \textbf{รายการประเมิน} & $\bar{x}$ & $S.D.$ & \textbf{ระดับ}\\
  \midrule
  \endfirsthead
  \caption[]{(ต่อ)}\\
  \toprule
  \textbf{รายการประเมิน} & $\bar{x}$ & $S.D.$ & \textbf{ระดับ}\\
  \midrule
  \endhead
  \bottomrule
  \endlastfoot
  1. ระบบตอบคำถามได้ตรงกับเนื้อหาของรายวิชา & 4.60 & 0.50 & มากที่สุด\\
  2. คำอธิบายเข้าใจง่าย & 4.53 & 0.57 & มากที่สุด\\
  3. ระบบช่วยให้แก้ไขข้อผิดพลาดของโค้ดได้เร็วขึ้น & 4.57 & 0.63 & มากที่สุด\\
  4. ระบบตอบสนองได้รวดเร็ว & 4.37 & 0.67 & มาก\\
  5. ระบบใช้งานง่าย & 4.63 & 0.49 & มากที่สุด\\
  \midrule
  \textbf{ภาพรวม} & \textbf{4.52} & \textbf{0.58} & \textbf{มากที่สุด}\\
\end{longtable}

% =====================================================================
%  บทที่ 5  สรุปผล อภิปรายผล และข้อเสนอแนะ
% =====================================================================
\chapter{สรุปผล อภิปรายผล และข้อเสนอแนะ}
\label{ch:conclusion}

\section{สรุปผลการวิจัย}

การวิจัยนี้ได้พัฒนาระบบผู้ช่วยสอนด้วยปัญญาประดิษฐ์เชิงสร้างสรรค์สำหรับรายวิชาการเขียนโปรแกรมเบื้องต้น ผลการทดลองพบว่าผู้เรียนที่ใช้ระบบมีผลสัมฤทธิ์ทางการเรียนสูงกว่าผู้เรียนที่ไม่ได้ใช้ระบบอย่างมีนัยสำคัญทางสถิติ (บทที่~\ref{ch:results}) และมีความพึงพอใจต่อระบบในระดับมากที่สุด

\section{อภิปรายผล}

ผลการวิจัยสอดคล้องกับแนวคิดที่ว่า การให้ข้อมูลป้อนกลับทันทีช่วยให้ผู้เรียนเข้าใจข้อผิดพลาดของตนเองได้ดีขึ้น นอกจากนี้ การใช้เทคนิค RAG ทำให้คำตอบของระบบอ้างอิงเนื้อหาของรายวิชา ลดปัญหาการสร้างข้อมูลที่ไม่ถูกต้อง (hallucination) ของแบบจำลองภาษา

\section{ข้อเสนอแนะ}

\subsection{ข้อเสนอแนะในการนำผลการวิจัยไปใช้}
\begin{enumerate}
  \item ผู้สอนควรกำหนดแนวทางการใช้งานระบบให้ชัดเจน เพื่อป้องกันการคัดลอกคำตอบโดยไม่ทำความเข้าใจ
  \item ควรปรับปรุงเอกสารประกอบการสอนในฐานข้อมูลให้เป็นปัจจุบันอยู่เสมอ
\end{enumerate}

\subsection{ข้อเสนอแนะในการวิจัยครั้งต่อไป}
\begin{enumerate}
  \item ควรศึกษาผลของระบบในระยะยาวตลอดภาคการศึกษา
  \item ควรขยายการใช้งานไปยังรายวิชาอื่น เช่น โครงสร้างข้อมูลและอัลกอริทึม
\end{enumerate}


% ---------- ส่วนท้าย ----------
\backmatter
\printbibliography[heading=thesis]

\appendix
\appendixpage                    % หน้าคั่น "ภาคผนวก"
% =====================================================================
%  ภาคผนวก ก
% =====================================================================
\chapter{แบบสอบถามความพึงพอใจ}
\label{app:questionnaire}

\noindent\textbf{คำชี้แจง} โปรดทำเครื่องหมาย \checkmark\ ในช่องที่ตรงกับความคิดเห็นของท่านมากที่สุด
(5 = มากที่สุด, 4 = มาก, 3 = ปานกลาง, 2 = น้อย, 1 = น้อยที่สุด)

\begin{table}[H]
  \caption{แบบสอบถามความพึงพอใจต่อระบบผู้ช่วยสอน}
  \label{tab:questionnaire}
  \centering
  \begin{tabular}{@{}p{0.6\linewidth}ccccc@{}}
    \toprule
    \textbf{รายการประเมิน} & \textbf{5} & \textbf{4} & \textbf{3} & \textbf{2} & \textbf{1}\\
    \midrule
    1. ระบบตอบคำถามได้ตรงกับเนื้อหาของรายวิชา & $\square$ & $\square$ & $\square$ & $\square$ & $\square$\\
    2. คำอธิบายเข้าใจง่าย & $\square$ & $\square$ & $\square$ & $\square$ & $\square$\\
    3. ระบบช่วยให้แก้ไขข้อผิดพลาดของโค้ดได้เร็วขึ้น & $\square$ & $\square$ & $\square$ & $\square$ & $\square$\\
    4. ระบบตอบสนองได้รวดเร็ว & $\square$ & $\square$ & $\square$ & $\square$ & $\square$\\
    5. ระบบใช้งานง่าย & $\square$ & $\square$ & $\square$ & $\square$ & $\square$\\
    \bottomrule
  \end{tabular}
\end{table}
    % ภาคผนวก ก
% =====================================================================
%  ภาคผนวก ข
% =====================================================================
\chapter{ตัวอย่างพรอมต์ที่ใช้ในระบบ}
\label{app:prompt}

\section{พรอมต์ระบบ (system prompt)}

\begin{lstlisting}[numbers=none]
You are a patient teaching assistant for an introductory Python course.
Answer using ONLY the course materials provided in <context>.
Explain the idea step by step and give a hint before giving code.
If the answer is not in the materials, say so.
\end{lstlisting}

\section{ตัวอย่างบทสนทนา}

\noindent\textbf{ผู้เรียน:} ทำไมโปรแกรมแจ้งว่า \texttt{IndexError: list index out of range}

\noindent\textbf{ระบบ:} ข้อผิดพลาดนี้เกิดขึ้นเมื่อเรียกใช้ตำแหน่งที่ไม่มีอยู่ในรายการ เช่น รายการที่มีสมาชิก 3 ตัวจะมีตำแหน่ง 0, 1 และ 2 เท่านั้น ลองตรวจสอบว่าเงื่อนไขของการวนซ้ำใช้ \texttt{range(len(items))} หรือไม่
    % ภาคผนวก ข

% =====================================================================
%  ประวัติผู้เขียน
%  รูปแบบ: \bioitem{หัวข้อ}{รายละเอียด}   (ใช้ \\ เพื่อขึ้นบรรทัดใหม่)
% =====================================================================
\begin{biography}
\bioitem{ชื่อ-นามสกุล}{นายสมชาย ใจดี}
\bioitem{วัน เดือน ปีเกิด}{1 มกราคม 2543}
\bioitem{ประวัติการศึกษา}{พ.ศ. 2565 วิศวกรรมศาสตรบัณฑิต สาขาวิชาวิศวกรรมคอมพิวเตอร์\\
  \theinstituteth}
\bioitem{ประสบการณ์การทำงาน}{พ.ศ. 2565--ปัจจุบัน วิศวกรซอฟต์แวร์ บริษัทตัวอย่าง จำกัด}
\bioitem{ผลงานตีพิมพ์}{สมชาย ใจดี และสมหญิง ตั้งใจเรียน. (2569). ชื่อบทความตัวอย่าง.
  \textit{ชื่อการประชุมวิชาการ}, 1--6.}
\end{biography}
     % ประวัติผู้เขียน

\end{document}
